\begin{proof}[Proof of \cref{obj:lem:fixed-budget-realization}]
\begin{enumerate}
\item \textbf{Admissibility of the dyadic-floor policy.}
For a real input \(v_{\mathrm{name}}\in\mathbb R\) and an integer precision \(q\ge0\), define
\[
m_q(v_{\mathrm{name}})
=\lfloor 2^qv_{\mathrm{name}}\rfloor,
\qquad
\mathcal J_q(v_{\mathrm{name}})
=
\left[
\frac{m_q(v_{\mathrm{name}})}{2^q},
\frac{m_q(v_{\mathrm{name}})+1}{2^q}
\right].
\]
The floor inequalities give
\[
m_q(v_{\mathrm{name}})
\le 2^qv_{\mathrm{name}}
<m_q(v_{\mathrm{name}})+1,
\]
and therefore
\[
v_{\mathrm{name}}\in\mathcal J_q(v_{\mathrm{name}}),
\qquad
\operatorname{wid}\bigl(\mathcal J_q(v_{\mathrm{name}})\bigr)=2^{-q},
\]
where throughout the proof
\[
\operatorname{wid}([u_{\mathrm{lo}},u_{\mathrm{hi}}])
=u_{\mathrm{hi}}-u_{\mathrm{lo}}.
\]
Moreover, multiplying the floor inequalities by two yields
\[
2m_q(v_{\mathrm{name}})
\le m_{q+1}(v_{\mathrm{name}})
<2\bigl(m_q(v_{\mathrm{name}})+1\bigr).
\]
Since the middle term is an integer, these inequalities imply
\[
\mathcal J_{q+1}(v_{\mathrm{name}})
\subseteq\mathcal J_q(v_{\mathrm{name}}).
\]
The conclusions include dyadic inputs, which may equal the lower endpoint of a queried interval. At each fixed precision, both rational endpoints are Borel functions of the real input because the floor function is Borel. Composing these functions with the public exponent coordinates and the observed covariate coordinates gives precisely the validity and Borel requirements in \cref{obj:def:dyadic-name-convention}. Thus \(\mathsf N_0\) is admissible.
% lean: CausalSmith.Stat.LogoddsLowsmoothFrontier.dyadicFloorPolicy_admissible

\item \textbf{Resolution enclosures at the rank budget.}
Fix an admissible engine \(\mathsf E\), an admissible policy \(\mathsf N\), \(n\ge2\), exponents in the stated domain, and a record tuple \(\mathcal O\). Set
\[
q_{\mathrm r}=q_{\mathrm{rank}}(n),
\qquad
\epsilon_{\mathrm r}=2^{-q_{\mathrm r}}.
\]
Let the exponent-name intervals after the intersections prescribed in \cref{obj:def:resolutions} be
\[
\mathcal J_{\alpha,\mathrm r}
=[\alpha_{\mathrm r,-},\alpha_{\mathrm r,+}],
\qquad
\mathcal J_{\beta,\mathrm r}
=[\beta_{\mathrm r,-},\beta_{\mathrm r,+}].
\]
Validity of the names and the exponent domain ensure that both intersections succeed and that
\[
\begin{gathered}
\alpha\in\mathcal J_{\alpha,\mathrm r}\subseteq[0,1],
\qquad
\beta\in\mathcal J_{\beta,\mathrm r}\subseteq[0,1/4],\\
\operatorname{wid}(\mathcal J_{\alpha,\mathrm r})\le\epsilon_{\mathrm r},
\qquad
\operatorname{wid}(\mathcal J_{\beta,\mathrm r})\le\epsilon_{\mathrm r}.
\end{gathered}
\]
Define the rational denominator and power-exponent boxes by exact interval operations:
\[
\begin{aligned}
\mathcal D_{C,\mathrm r}
&=2\mathcal J_{\alpha,\mathrm r}
  +2\mathcal J_{\beta,\mathrm r}+[1,1],
&
\mathcal D_{S,\mathrm r}
&=4\mathcal J_{\beta,\mathrm r}+[1,1],\\
\mathcal P_{C,\mathrm r}
&=[2,2]/\mathcal D_{C,\mathrm r},
&
\mathcal P_{S,\mathrm r}
&=[2,2]/\mathcal D_{S,\mathrm r}.
\end{aligned}
\]
Both denominator boxes have lower endpoints at least one and widths at most \(4\epsilon_{\mathrm r}\). For any denominator box
\(\mathcal D_{\mathrm{aux}}=[d_{\mathrm{lo}},d_{\mathrm{hi}}]\) with
\(1\le d_{\mathrm{lo}}\le d_{\mathrm{hi}}\), exact division gives
\[
[2,2]/\mathcal D_{\mathrm{aux}}
=
[2/d_{\mathrm{hi}},2/d_{\mathrm{lo}}],
\]
and
\[
\operatorname{wid}([2,2]/\mathcal D_{\mathrm{aux}})
=
\frac{2(d_{\mathrm{hi}}-d_{\mathrm{lo}})}
     {d_{\mathrm{lo}}d_{\mathrm{hi}}}
\le2\operatorname{wid}(\mathcal D_{\mathrm{aux}}).
\]
Consequently, both power-exponent boxes lie in \([0,2]\), contain their respective exact exponents, and have width at most \(8\epsilon_{\mathrm r}\).

For a power exponent \(\tau\in[0,2]\),
\[
\left|\frac{d}{d\tau}n^\tau\right|
=n^\tau\log n
\le n^2n=n^3,
\]
using \(n\ge1\) and \(\log n\le n\). The mean-value bound and the outward endpoint errors required by \cref{obj:def:arithmetic-engine} therefore show that the returned target boxes satisfy
\[
R_C\in[a_C,b_C],\qquad R_S\in[a_S,b_S],
\]
\[
b_C-a_C,\ b_S-a_S
\le8n^3\epsilon_{\mathrm r}+2\epsilon_{\mathrm r}.
\]
To pay for this width, the integer logarithmic ceiling in \cref{obj:def:resolutions} gives
\[
n\le2^{\lceil\log_2n\rceil},
\qquad
128n^3\epsilon_{\mathrm r}
=
\frac{n^3}{2^{3\lceil\log_2n\rceil}}
\le1.
\]
Since \(n^3\ge1\), it follows that
\[
8n^3\epsilon_{\mathrm r}+2\epsilon_{\mathrm r}
\le10n^3\epsilon_{\mathrm r}
\le\frac{10}{128}\le1.
\]
All rank-step intersections and divisions have thus succeeded, the power inputs satisfy the primitive domain requirements, and both returned boxes have the required unit-width bound.
% lean: CausalSmith.Stat.LogoddsLowsmoothFrontier.rankBoxes_rank_budget_contract

\item \textbf{Bounds for the selected ranks.}
The exponent domain makes both target exponents positive. Hence
\[
R_C\ge1,\qquad R_S\ge1.
\]
In particular \(b_C,b_S\ge1\), so the nonnegative ceilings in \cref{obj:def:resolutions} are their ordinary integer ceilings. Containment and the unit-width bounds give
\[
R_C\le b_C\le a_C+1\le R_C+1,
\qquad
R_S\le b_S\le a_S+1\le R_S+1.
\]
Using the strict upper bound for an integer ceiling,
\[
\begin{aligned}
R_C&\le k_C=\lceil b_C\rceil<b_C+1\le R_C+2\le3R_C,\\
R_S&\le k_S=\lceil b_S\rceil<b_S+1\le R_S+2\le3R_S.
\end{aligned}
\]
The last inequalities follow from \(R_C,R_S\ge1\). These ranks depend on the public exponent names and are fixed before sampling. In particular,
\[
k_C\ge1,\qquad k_S\ge1.
\]
From this point onward, \(b_C,b_S\) denote the coordinate error radii in \cref{obj:def:upper-handle}.
% lean: CausalSmith.Stat.LogoddsLowsmoothFrontier.enclosure_ceiling_rank_bounds

\item \textbf{Coordinate enclosures at the endpoint budget.}
Retain the selected ranks and set
\[
K=\max\{k_C,k_S\},
\qquad
q_{\mathrm e}=q_{\mathrm{end}}(n,k_C,k_S),
\qquad
\epsilon_{\mathrm e}=2^{-q_{\mathrm e}}.
\]
At this precision, write the raw name responses as
\[
\begin{gathered}
\mathcal B_\alpha
=[\alpha_{\mathrm{name},-},\alpha_{\mathrm{name},+}],
\qquad
\mathcal B_\beta
=[\beta_{\mathrm{name},-},\beta_{\mathrm{name},+}],\\
\mathcal B_{X_i}
=[X_{i,\mathrm{name},-},X_{i,\mathrm{name},+}],
\qquad 1\le i\le n.
\end{gathered}
\]
These boxes contain their named inputs and have widths at most \(\epsilon_{\mathrm e}\). Since \(q_{\mathrm e}\ge2\), one has
\[
\epsilon_{\mathrm e}\le\frac14.
\]
Containment and the width bound place each raw covariate box in
\[
\mathcal B_{X_i}\subseteq[-\epsilon_{\mathrm e},1+\epsilon_{\mathrm e}].
\]
Indeed, each endpoint is within one box width of the contained covariate.

For any displayed scalar expression \(G\), denote its computed enclosure in the prescribed tree at precision \(q_{\mathrm e}\) by
\[
\mathcal B[G]=[G_{\mathrm{box},-},G_{\mathrm{box},+}].
\]
This notation refers to the supplied raw responses, the fixed engine, and the exact rational range operations.

We use two elementary properties of these operations. For closed rational boxes \(\mathcal J_1,\mathcal J_2\), define
\[
M_{\mathcal J_\iota}
=\max\{|u|:u\in\mathcal J_\iota\},
\qquad \iota\in\{1,2\}.
\]
Inserting a cross product when subtracting two products, and then taking the extrema among endpoint products, gives
\[
\operatorname{wid}(\mathcal J_1\mathcal J_2)
\le
M_{\mathcal J_2}\operatorname{wid}(\mathcal J_1)
+
M_{\mathcal J_1}\operatorname{wid}(\mathcal J_2).
\]
Widths add under addition and subtraction. The endpoint formulas for minimum and maximum also give
\[
\begin{aligned}
\operatorname{wid}(\min(\mathcal J_1,\mathcal J_2))
&\le\max\{\operatorname{wid}(\mathcal J_1),
          \operatorname{wid}(\mathcal J_2)\},\\
\operatorname{wid}(\max(\mathcal J_1,\mathcal J_2))
&\le\max\{\operatorname{wid}(\mathcal J_1),
          \operatorname{wid}(\mathcal J_2)\}.
\end{aligned}
\]
For example, each upper endpoint is at most its corresponding lower endpoint plus the larger width; taking minima or maxima preserves that inequality.

We now evaluate the basis and kernel from \cref{obj:synth_3}. Define the primitive constant boxes
\[
\mathcal B_\pi=\mathcal B[\pi],
\qquad
\mathcal B_{\sqrt2}=\mathcal B[\sqrt2].
\]
Their contracts give
\[
\mathcal B_\pi\subseteq[3,4],
\qquad
\operatorname{wid}(\mathcal B_\pi)\le2\epsilon_{\mathrm e},
\]
\[
\mathcal B_{\sqrt2}\subseteq[0,2],
\qquad
\operatorname{wid}(\mathcal B_{\sqrt2})\le2\epsilon_{\mathrm e}.
\]
For the second inclusion, \(\sqrt2\le3/2\) follows by squaring, and the upper endpoint error is at most \(\epsilon_{\mathrm e}\le1/4\).

For a basis index \(j\ge1\), the intermediate argument boxes are
\[
\mathcal B_{\pi j}=\mathcal B_\pi[j,j],
\qquad
\mathcal B_{\pi jX_i}=\mathcal B_{\pi j}\mathcal B_{X_i}.
\]
They satisfy
\[
\operatorname{wid}(\mathcal B_{\pi j})\le2j\epsilon_{\mathrm e},
\qquad
M_{\mathcal B_{\pi j}}\le4j,
\qquad
M_{\mathcal B_{X_i}}\le\frac54.
\]
Thus the product-width inequality yields
\[
\operatorname{wid}(\mathcal B_{\pi jX_i})
\le
\frac54(2j\epsilon_{\mathrm e})+4j\epsilon_{\mathrm e}
=\frac{13}{2}j\epsilon_{\mathrm e}.
\]
The cosine contract returns a box in \([-1,1]\) whose width is at most
\[
\operatorname{wid}\bigl(\mathcal B[\cos(\pi jX_i)]\bigr)
\le\left(\frac{13}{2}j+2\right)\epsilon_{\mathrm e}.
\]
Multiplication by the square-root box therefore gives
\[
M_{\mathcal B[\varphi_j(X_i)]}\le2,
\]
\[
\operatorname{wid}\bigl(\mathcal B[\varphi_j(X_i)]\bigr)
\le
2\epsilon_{\mathrm e}
+2\left(\frac{13}{2}j+2\right)\epsilon_{\mathrm e}
=(13j+6)\epsilon_{\mathrm e}
\le20(j+1)\epsilon_{\mathrm e}.
\]
For \(j=0\), the basis expression is the exact constant one, so these magnitude and width bounds also hold.

Each computed basis-product term consequently has width at most
\(80(j+1)\epsilon_{\mathrm e}\). Summing the terms in the kernel, for either retained rank \(k\in\{k_C,k_S\}\), gives
\[
\begin{aligned}
\operatorname{wid}\bigl(\mathcal B[H_k(X_i,X_{i'})]\bigr)
&\le\sum_{j=0}^{k-1}80(j+1)\epsilon_{\mathrm e}\\
&\le\sum_{j=0}^{k-1}80k\epsilon_{\mathrm e}
=80k^2\epsilon_{\mathrm e}
\le160k^2\epsilon_{\mathrm e}.
\end{aligned}
\]
At \(k=1\), the sole basis-product term is exact.

The binary marks used in \cref{obj:def:coordinate-estimates} are exact rational numbers in \([0,1]\). Multiplication by their products preserves the displayed kernel-width bound. There are exactly
\[
\#\{(i,i'):1\le i,i'\le n,\ i\ne i'\}=n(n-1)
\]
ordered pairs. Hence the positive rational normalization in \cref{obj:def:projection-statistic} cancels this factor in the sum of widths:
\[
\operatorname{wid}\bigl(\mathcal B[\mathbb U_{n,k}(W,V)]\bigr)
\le
\frac{n(n-1)\,160k^2\epsilon_{\mathrm e}}{n(n-1)}
=160k^2\epsilon_{\mathrm e}
\]
for the marks used here. The empirical mean in \(\widehat C_n\) is exact. Soundness of the operations therefore supplies successful coordinate boxes with
\[
\begin{gathered}
\widehat C_n\in\mathcal B[\widehat C_n],
\qquad
\widehat S_n\in\mathcal B[\widehat S_n],\\
\operatorname{wid}\bigl(\mathcal B[\widehat C_n]\bigr)
\le160k_C^2\epsilon_{\mathrm e},
\qquad
\operatorname{wid}\bigl(\mathcal B[\widehat S_n]\bigr)
\le160k_S^2\epsilon_{\mathrm e}.
\end{gathered}
\]
% lean: CausalSmith.Stat.LogoddsLowsmoothFrontier.endpoint_coordinates_budget_contract

\item \textbf{Error-radius enclosures.}
The endpoint budget gives
\[
K\le2^{\lceil\log_2(K+1)\rceil},
\qquad
q_{\mathrm e}\ge4+\lceil\log_2(K+1)\rceil,
\]
and therefore
\[
K\epsilon_{\mathrm e}\le\frac1{16}.
\]
Validity of the raw exponent responses gives
\[
\mathcal B_\alpha\subseteq[-\epsilon_{\mathrm e},1+\epsilon_{\mathrm e}],
\qquad
\mathcal B_\beta\subseteq[-\epsilon_{\mathrm e},1/4+\epsilon_{\mathrm e}].
\]
Define their bias-exponent boxes by
\[
\mathcal B_{C,\mathrm{pow}}=[0,0]-(\mathcal B_\alpha+\mathcal B_\beta),
\qquad
\mathcal B_{S,\mathrm{pow}}=[-2,-2]\mathcal B_\beta.
\]
They contain \(-(\alpha+\beta)\) and \(-2\beta\), respectively, and satisfy
\[
\begin{gathered}
\mathcal B_{C,\mathrm{pow}}
\subseteq[-5/4-2\epsilon_{\mathrm e},2\epsilon_{\mathrm e}]
\subseteq[-3,3],\\
\mathcal B_{S,\mathrm{pow}}
\subseteq[-1/2-2\epsilon_{\mathrm e},2\epsilon_{\mathrm e}]
\subseteq[-3,3],\\
\operatorname{wid}(\mathcal B_{C,\mathrm{pow}}),
\ \operatorname{wid}(\mathcal B_{S,\mathrm{pow}})
\le2\epsilon_{\mathrm e}.
\end{gathered}
\]

For either retained rank \(k\), the elementary inequality
\(\log u\le u-1\) for \(u>0\) gives
\[
\log k=\log(k/2)+\log2
\le(k/2-1)+1=k/2.
\]
Applying the same inequality at \(u=1/2\) gives \(\log2\ge1/2\).
Thus, for every exponent \(\tau\) in the corresponding raw power-exponent box,
\[
\tau\log k
\le2\epsilon_{\mathrm e}\log k
\le k\epsilon_{\mathrm e}
\le\frac1{16}\le\log2.
\]
It follows that
\[
k^\tau\le2,
\qquad
\left|\frac{d}{d\tau}k^\tau\right|
=k^\tau\log k\le k.
\]
These inequalities include \(k=1\), for which the derivative is zero.
The mean-value bound and the primitive endpoint errors now imply
\[
\operatorname{wid}\bigl(\mathcal B[k^\tau]\bigr)
\le k\operatorname{wid}(\mathcal B_{\mathrm{pow}})
    +2\epsilon_{\mathrm e}
\le(2k+2)\epsilon_{\mathrm e},
\]
with \(\mathcal B_{\mathrm{pow}}\) the corresponding bias-exponent box.

The variance quantity
\[
V_n(k)=\frac{10}{n}+\frac{4k}{n(n-1)}
\]
is an exactly evaluated positive rational number. Its square-root enclosure therefore satisfies
\[
\operatorname{wid}\bigl(\mathcal B[\sqrt{20V_n(k)}]\bigr)
\le2\epsilon_{\mathrm e}.
\]
Using the constants in the radius expressions, we obtain successful enclosures containing \(b_C,b_S\), with
\[
\begin{aligned}
\operatorname{wid}(\mathcal B[b_C])
&\le\bigl(8(2k_C+2)+2\bigr)\epsilon_{\mathrm e}
=(16k_C+18)\epsilon_{\mathrm e}
\le34K\epsilon_{\mathrm e},\\
\operatorname{wid}(\mathcal B[b_S])
&\le\bigl(25(2k_S+2)+2\bigr)\epsilon_{\mathrm e}
=(50k_S+52)\epsilon_{\mathrm e}
\le102K\epsilon_{\mathrm e}.
\end{aligned}
\]
The last inequalities use \(k_C,k_S\le K\) and \(K\ge1\).
% lean: CausalSmith.Stat.LogoddsLowsmoothFrontier.endpoint_radii_budget_contract

\item \textbf{Denominator-corner enclosures.}
Define the computed denominator boxes by the prescribed exact clipping tree:
\[
\mathcal B[s_\pm]
=
\max\left(
[s_0,s_0],
\min\left([1,1],
\mathcal B[\widehat S_n]\pm\mathcal B[b_S]\right)
\right),
\qquad s_0=\frac1{512}.
\]
Here the minus sign denotes interval subtraction. Soundness supplies containment of \(s_\pm\), and the endpoint formulas give
\[
\mathcal B[s_\pm]\subseteq[1/512,1].
\]
Before clipping, the widths are bounded by
\[
160k_S^2\epsilon_{\mathrm e}+102K\epsilon_{\mathrm e}
\le262K^2\epsilon_{\mathrm e}
\le300K^2\epsilon_{\mathrm e},
\]
because \(k_S\le K\) and \(K\le K^2\). Clipping by point intervals preserves this bound by the minimum and maximum width inequalities. Consequently,
\[
s_\pm\in\mathcal B[s_\pm],
\qquad
\operatorname{wid}(\mathcal B[s_\pm])
\le300K^2\epsilon_{\mathrm e}.
\]
In particular, every subsequent denominator interval has a positive rational lower endpoint.
% lean: CausalSmith.Stat.LogoddsLowsmoothFrontier.endpoint_denominator_corners

\item \textbf{Numerator-corner widths and magnitudes.}
Define
\[
\mathcal B[c_\pm]
=\mathcal B[\widehat C_n]\pm\mathcal B[b_C].
\]
These boxes contain \(c_\pm\), and
\[
\operatorname{wid}(\mathcal B[c_\pm])
\le160k_C^2\epsilon_{\mathrm e}+34K\epsilon_{\mathrm e}
\le194K^2\epsilon_{\mathrm e}
\le300K^2\epsilon_{\mathrm e}.
\]

We also require bounds on their endpoints. The exact basis satisfies
\[
|\varphi_j(x)|\le\sqrt2
\qquad(j\ge0,\ x\in[0,1]),
\]
including \(j=0\). Therefore
\[
|H_k(x,z)|
\le\sum_{j=0}^{k-1}|\varphi_j(x)\varphi_j(z)|
\le2k.
\]
The ordered-pair normalization and marks of absolute value at most one give
\[
|\mathbb U_{n,k}(W,V)|\le2k.
\]
The empirical treated-outcome mean has absolute value at most one, so
\[
|\widehat C_n|\le1+2k_C.
\]
Since \(n\ge2\),
\[
V_n(k)\le5+2k,
\qquad
20V_n(k)\le100+40k\le(10+7k)^2.
\]
Taking square roots yields
\[
\sqrt{20V_n(k)}\le10+7k.
\]
Moreover, \(k_C\ge1\) and \(\alpha+\beta>0\) give
\[
k_C^{-(\alpha+\beta)}\le1,
\qquad
0\le b_C\le8+10+7k_C\le25k_C.
\]
Thus
\[
|c_\pm|
\le|\widehat C_n|+b_C
\le1+27k_C\le1+27K.
\]

To control the computed endpoint overhang, the two logarithmic ceilings in the endpoint budget give
\[
n+1\le2^{\lceil\log_2(n+1)\rceil},
\qquad
K+1\le2^{\lceil\log_2(K+1)\rceil},
\]
hence
\[
2^{44}(n+1)(K+1)^4\epsilon_{\mathrm e}\le1.
\]
Comparing coefficients with this product shows
\[
\epsilon_{\mathrm e}\le\frac1{16},
\qquad
300K^2\epsilon_{\mathrm e}\le1,
\qquad
2^{37}(K+1)^3\epsilon_{\mathrm e}\le\frac1n.
\]
For the last two comparisons, use
\[
K^2\le(K+1)^4,
\qquad
300K^2\le2^{44}(n+1)(K+1)^4,
\]
\[
2^{37}n(K+1)^3
\le2^{44}(n+1)(K+1)^4.
\]

A box containing \(c_\pm\) and having width at most one has every point within distance one of \(c_\pm\). It follows that
\[
M_{\mathcal B[c_\pm]}
\le1+27K+1\le100(K+1).
\]
This supplies both the width and magnitude bounds at the actual numerator inputs to inversion.
% lean: CausalSmith.Stat.LogoddsLowsmoothFrontier.endpoint_numerator_corners

\item \textbf{Enclosures of the four inversion corners.}
For signs \(\sigma_c,\sigma_s\in\{-,+\}\), define
\[
\mathcal B_{\mathrm{quot},\sigma_c\sigma_s}
=\mathcal B[c_{\sigma_c}]/\mathcal B[s_{\sigma_s}],
\qquad
\mathcal B_{\mathrm{odds},\sigma_c\sigma_s}
=[1,1]+\mathcal B_{\mathrm{quot},\sigma_c\sigma_s}.
\]
For a denominator box
\(\mathcal J_{\mathrm{den}}=[d_{\mathrm{lo}},d_{\mathrm{hi}}]\)
with \(d_{\mathrm{lo}}\ge1/512\), its reciprocal range satisfies
\[
\mathcal J_{\mathrm{den}}^{-1}
=[1/d_{\mathrm{hi}},1/d_{\mathrm{lo}}],
\qquad
M_{\mathcal J_{\mathrm{den}}^{-1}}\le512,
\]
\[
\operatorname{wid}(\mathcal J_{\mathrm{den}}^{-1})
=
\frac{d_{\mathrm{hi}}-d_{\mathrm{lo}}}
     {d_{\mathrm{lo}}d_{\mathrm{hi}}}
\le512^2\operatorname{wid}(\mathcal J_{\mathrm{den}}).
\]
Exact division is multiplication by this reciprocal range. The product-width bound therefore gives containment of the true ratio and
\[
\begin{aligned}
\operatorname{wid}(\mathcal B_{\mathrm{quot},\sigma_c\sigma_s})
\le{}&
512\operatorname{wid}(\mathcal B[c_{\sigma_c}])\\
&+100(K+1)512^2
 \operatorname{wid}(\mathcal B[s_{\sigma_s}]).
\end{aligned}
\]

Define the exponential boundary boxes and the clipped odds boxes by
\[
\mathcal B_{\exp,-}=\mathcal B[\exp(-1/2)],
\qquad
\mathcal B_{\exp,+}=\mathcal B[\exp(1/2)],
\]
\[
\mathcal B_{\mathrm{clip},\sigma_c\sigma_s}
=
\max\left(
\mathcal B_{\exp,-},
\min\left(\mathcal B_{\exp,+},
          \mathcal B_{\mathrm{odds},\sigma_c\sigma_s}\right)
\right).
\]
The exponential series, bounded termwise by the geometric series, gives
\[
\exp(1/4)\le\sum_{j=0}^{\infty}(1/4)^j=\frac43,
\]
and consequently
\[
\exp(1/2)=\exp(1/4)^2\le\frac{16}{9},
\qquad
\exp(-1/2)\ge\frac9{16}.
\]
The primitive endpoint errors and \(\epsilon_{\mathrm e}\le1/16\) imply
\[
\inf\mathcal B_{\exp,-}\ge\frac9{16}-\epsilon_{\mathrm e}\ge\frac12,
\qquad
\sup\mathcal B_{\exp,+}\le\frac{16}{9}+\epsilon_{\mathrm e}\le2,
\]
\[
\operatorname{wid}(\mathcal B_{\exp,-}),
\ \operatorname{wid}(\mathcal B_{\exp,+})
\le2\epsilon_{\mathrm e}.
\]
Thus every clipped odds box has lower endpoint at least \(1/2\), contains the exact clipped odds, and satisfies
\[
\operatorname{wid}(\mathcal B_{\mathrm{clip},\sigma_c\sigma_s})
\le
\operatorname{wid}(\mathcal B_{\mathrm{quot},\sigma_c\sigma_s})
+2\epsilon_{\mathrm e}.
\]
Here the added exact constant one preserves width, and the minimum and maximum width inequalities bound both clipping operations.

Define the inversion-corner boxes by
\[
\mathcal B_{F,\sigma_c\sigma_s}
=\mathcal B[F(c_{\sigma_c},s_{\sigma_s})].
\]
The logarithm derivative is at most two on every argument interval with lower endpoint at least \(1/2\). Its primitive contract therefore supplies successful boxes containing the four exact corner values, with
\[
\begin{aligned}
\operatorname{wid}(\mathcal B_{F,\sigma_c\sigma_s})
\le{}&
2\left(
512\operatorname{wid}(\mathcal B[c_{\sigma_c}])
+100(K+1)512^2
 \operatorname{wid}(\mathcal B[s_{\sigma_s}])
\right)
+6\epsilon_{\mathrm e}\\
\le{}&
\bigl(307200K^2
+15728640000(K+1)K^2+6\bigr)\epsilon_{\mathrm e}\\
\le{}&2^{37}(K+1)^3\epsilon_{\mathrm e}
\le\frac1n.
\end{aligned}
\]
For the coefficient comparison, \(K\ge0\) implies
\[
K^2\le(K+1)^3,\qquad
(K+1)K^2\le(K+1)^3,\qquad
1\le(K+1)^3,
\]
and
\[
307200+15728640000+6\le2^{37}.
\]
This establishes the required precision for each of the four corner evaluations.
% lean: CausalSmith.Stat.LogoddsLowsmoothFrontier.endpoint_inversion_budget_contract

\item \textbf{The two endpoint boxes.}
The literal minimum and maximum operations return
\[
\mathcal B_L
=\min(\mathcal B_{F,--},\mathcal B_{F,-+})
=[L_-,L_+],
\]
\[
\mathcal B_R
=\max(\mathcal B_{F,+-},\mathcal B_{F,++})
=[R_-,R_+].
\]
Monotonicity of minimum and maximum in their arguments preserves containment, giving
\[
L\in[L_-,L_+],
\qquad
R\in[R_-,R_+].
\]
Their widths satisfy
\[
L_+-L_-
\le\max\{
\operatorname{wid}(\mathcal B_{F,--}),
\operatorname{wid}(\mathcal B_{F,-+})\}
\le\frac1n,
\]
\[
R_+-R_-
\le\max\{
\operatorname{wid}(\mathcal B_{F,+-}),
\operatorname{wid}(\mathcal B_{F,++})\}
\le\frac1n.
\]
All raw input queries and all intervening evaluations have succeeded. Every power input lies in its required domain, every square-root input is nonnegative, every division has a positive denominator interval, and every logarithm input has a positive lower endpoint. The finite sums at the retained finite ranks and the primitive termination contracts in \cref{obj:def:arithmetic-engine} therefore give successful return of both boxes at the single endpoint precision.
% lean: CausalSmith.Stat.LogoddsLowsmoothFrontier.endpoint_min_max_width

\item \textbf{Outward enclosure and rational interval shape.}
The exact clipped logarithmic map in \cref{obj:def:upper-handle} satisfies
\[
-\frac12\le F(c,d)\le\frac12
\]
at the endpoint inputs: its logarithm argument lies between
\(\exp(-1/2)\) and \(\exp(1/2)\). For each positive denominator, division, clipping, and logarithm are nondecreasing in the numerator. Since \(b_C\ge0\) and \(s_-\ge s_0>0\),
\[
L
\le F(c_-,s_-)
\le F(c_+,s_-)
\le R.
\]
Consequently,
\[
-\frac12\le L\le R\le\frac12.
\]
Define the literal rational output endpoints by
\[
l_{\mathrm{out}}=\max\{-1/2,L_-\},
\qquad
u_{\mathrm{out}}=\min\{1/2,R_+\}.
\]
Containment of the exact endpoints and their effect bounds give
\[
l_{\mathrm{out}}\le L\le R\le u_{\mathrm{out}}.
\]
Thus the intersection in the output construction succeeds and the literal output is
\[
I^{\mathrm{up}}_{n,\mathsf E,\mathsf N}(\alpha,\beta)(\mathcal O)
=[l_{\mathrm{out}},u_{\mathrm{out}}],
\qquad
[L,R]\subseteq[l_{\mathrm{out}},u_{\mathrm{out}}].
\]
The endpoint-width bounds give
\[
0\le L-l_{\mathrm{out}}
\le L-L_-\le L_+-L_-\le\frac1n,
\]
\[
0\le u_{\mathrm{out}}-R
\le R_+-R\le R_+-R_-\le\frac1n.
\]
Adding these inequalities yields
\[
u_{\mathrm{out}}-l_{\mathrm{out}}
\le R-L+\frac2n.
\]
Finally,
\[
l_{\mathrm{out}},u_{\mathrm{out}}\in\mathbb Q,
\qquad
-\frac12\le l_{\mathrm{out}}\le u_{\mathrm{out}}\le\frac12.
\]
The output is therefore a nonempty closed connected interval with the asserted rational endpoint shape.
% lean: CausalSmith.Stat.LogoddsLowsmoothFrontier.upperOutput_enclosure_guarantees

\item \textbf{Borel measurability.}
Keep the engine, policy, sample size, and public exponents fixed. The ranks and endpoint precision are then fixed integers. Define the finite response code and its countable range by
\[
\mathfrak z(\mathcal O)
=
\left(
(A_i,Y_i),
(X_{i,\mathrm{name},-},X_{i,\mathrm{name},+})
\right)_{i=1}^n,
\qquad
\mathcal Z
=
\bigl(\{0,1\}^2\times\mathbb Q^2\bigr)^n.
\]
The binary coordinates and the fixed-precision name responses are Borel, so \(\mathfrak z\) is Borel into the discrete countable space \(\mathcal Z\). The public exponent responses are fixed, and the remaining computation uses this code through finite rational operations. Hence
\[
\mathfrak z(\mathcal O)=\mathfrak z(\mathcal O')
\quad\Longrightarrow\quad
I^{\mathrm{up}}_{n,\mathsf E,\mathsf N}(\alpha,\beta)(\mathcal O)
=
I^{\mathrm{up}}_{n,\mathsf E,\mathsf N}(\alpha,\beta)(\mathcal O').
\]
For any set \(\mathcal T\) of output intervals, define
\[
\mathcal Z_{\mathcal T}
=
\left\{
z\in\mathcal Z:
\exists\mathcal O,\ 
\mathfrak z(\mathcal O)=z,\
I^{\mathrm{up}}_{n,\mathsf E,\mathsf N}(\alpha,\beta)(\mathcal O)
\in\mathcal T
\right\}.
\]
The preceding implication gives the fiber decomposition
\[
\left\{
\mathcal O:
I^{\mathrm{up}}_{n,\mathsf E,\mathsf N}(\alpha,\beta)(\mathcal O)
\in\mathcal T
\right\}
=
\bigcup_{z\in\mathcal Z_{\mathcal T}}
\mathfrak z^{-1}(\{z\}).
\]
This is a countable union of Borel sets. In particular, inverse images of Borel sets of intervals are Borel, establishing the claimed measurability of the output map.
% lean: CausalSmith.Stat.LogoddsLowsmoothFrontier.upperOutput_regular
\end{enumerate}
\end{proof}
